\documentclass[11pt]{article}
\usepackage[margin=1in]{geometry}
\usepackage{graphicx,booktabs,float,xcolor,hyperref,amssymb}
\usepackage{helvet}\renewcommand{\familydefault}{\sfdefault}
\hypersetup{colorlinks=true,linkcolor=blue,urlcolor=blue}
\newcommand{\Mw}{$M_{\mathrm w}$}
\setlength{\parindent}{0pt}\setlength{\parskip}{6pt}
\title{\textbf{Prospective source model of the 2026 \Mw~7.4 San Jos\'e del Palmar (Colombia)
intermediate-depth earthquake}\\
\large A blind test of the MTUQ$+$WISP pipeline, frozen before the USGS finite fault ---
and the discovery of a 35-s emergent nucleation phase}
\author{Automated pipeline (17.Venezuela\_2026)}
\date{Frozen 2026-08-11 --- no USGS finite-fault product existed at freeze time}
\begin{document}\maketitle

\section*{Purpose and headline}
The \Mw~7.4 earthquake of 2026-08-10 beneath western Colombia (110 km deep) has, at freeze
time, no published finite-fault model. We freeze a teleseismic-only source model for later
head-to-head comparison. Two findings lead: (1) the event is \textbf{not} a Bucaramanga-nest
earthquake --- it belongs to the Cauca intermediate-depth cluster of the subducting Nazca slab,
407 km southwest of the nest; (2) the rupture began with a \textbf{$\sim$35-s weak emergent
nucleation phase} before the main moment release 29 km south of the epicentre --- a property
that breaks naive origin-anchored inversions and had to be handled explicitly.

\section{Event, tectonic setting, and the Bucaramanga question}
\begin{table}[H]\centering\small
\begin{tabular}{ll}
\toprule
Origin & 2026-08-10T12:34:28.1 UTC, 4.844$^\circ$N, 76.242$^\circ$W, depth 110.3 km \\
Magnitude & \Mw~7.4 (W-phase); $M_0=1.66\times10^{20}$ N\,m; centroid 120.5 km; DC 87\,\% \\
W-phase NP1 & 124.0/58.6/158.5 \\
W-phase NP2 & 225.6/71.8/33.2 \quad (NE--SW, steep, oblique reverse) \\
W-phase centroid & 4.585$^\circ$N, 76.161$^\circ$W (29 km S of epicentre), \textbf{12:35:02.6 = origin $+$ 34.5 s} \\
\bottomrule
\end{tabular}
\end{table}

\textbf{Bucaramanga nest? No.} The nest --- the world's densest intermediate-depth seismic
nest --- sits at $\sim$6.82$^\circ$N, 73.15$^\circ$W, 140--160 km depth beneath Santander, NE
Colombia. This event is \textbf{407 km to the southwest} at 4.84$^\circ$N, 76.24$^\circ$W,
110--120 km deep: the \textbf{Cauca intermediate-depth cluster} of the subducting Nazca slab
(cf.\ Chang et al.\ 2019), a different slab segment with its own seismicity. It ranks among
the largest instrumental events of that segment.

\section{Data}
120 three-component stations (316 channels), 30--90$^\circ$, maximum azimuthal gap 47$^\circ$
(the Pacific). MTUQ used an azimuth-balanced 36-station response-corrected subset with the
deep-event bands (body 25--60 s, mantle surface waves 200--450 s). Both nodal planes extend
below the 126.5-km WISP surface-wave GF cap, so the finite-fault plane comparison is
\textbf{body-only on both planes --- fair by construction} (the San Pedro de Atacama rule).

\begin{figure}[H]\centering
\includegraphics[width=0.6\textwidth]{../fig_colombia_stations.png}
\caption{Teleseismic stations (azimuthal-equidistant; rings 30/60/90$^\circ$).}
\end{figure}

\section{The emergent nucleation phase --- and why origin-anchored inversions fail}
The observed P wavetrains begin weakly and emergently; the main impulsive packet arrives only
$\sim$45--55 s after origin (its pP/sP depth-phase echoes at $\sim$75--90 s). Independently,
the W-phase centroid time lags the origin by 34.5 s and sits 29 km south. The implied
nucleation-to-asperity apparent velocity is $\sim$0.8 km\,s$^{-1}$ --- far below any
physical rupture speed --- so a finite fault nucleating at the cataloged hypocentre at the
cataloged time \emph{cannot} place the moment where the data demand it. Our first
(origin-anchored) inversion demonstrated the failure exactly: slip piled at the fault edge
(86\,\% of peak in an edge row --- a violated containment rule), the STF peaked at an
unreachable 40 s, and the synthetics fit only the long-period envelope. Origin-cut MTUQ
windows likewise missed the main energy (\Mw\ pinned at the grid edge).

\textbf{The fix} (the operational Yeck-style relocation): re-anchor the entire analysis ---
window cutting, Green's functions, fault geometry, nucleation --- at the \textbf{W-phase
centroid in space and time}. The 35-s precursor then lies outside the modelling window (its
moment is a few per cent at most), and the main rupture is imaged about its true origin.

\begin{figure}[H]\centering
\includegraphics[width=0.55\textwidth]{../20260810123428/ffm.0/NP1_shift/plots/MomentRate.png}
\caption{The INVALID origin-anchored model's STF, shown as the diagnostic: the peak at
$\sim$40 s is kinematically unreachable from the hypocentre, so the inversion smeared slip to
the fault edges.}
\end{figure}

\section{Finite fault (WISP, centroid-anchored): both planes, \texttt{shift\_match}, body-only}
\begin{table}[H]\centering\small
\begin{tabular}{llll}
\toprule
plane (W-phase labels) & misfit (post-shift) & edge slip top/bottom & peak slip \\
\midrule
\textbf{226/72/33 (preferred)} & \textbf{0.1363} & 1\,\% / 3\,\% & 1.63 m \\
124/59/159 conjugate & 0.1485 & \textbf{40\,\%} / 2\,\% & 1.12 m \\
\bottomrule
\end{tabular}
\caption{An \textbf{8.2\,\% margin plus the conjugate's top-edge violation = a near-decisive
plane call} for the NE--SW steep oblique-reverse plane --- exactly where the rupture-size
criterion (Calama 1.6\,\% tie at 10 km; SPdA 10\,\% decisive at 50 km) predicts discrimination
turns on. WISP's directory labels are swapped vs USGS (WISP \texttt{NP1} = W-phase NP2), as at
Kumamoto. Recovered $M_0=1.68\times10^{20}$ N\,m (W-phase 1.66).}
\end{table}

\begin{figure}[H]\centering
\includegraphics[width=0.85\textwidth]{../20260810123502/ffm.0/NP1_shift/plots/SlipIsochrone_trueScale.png}
\caption{Preferred plane (226/72/33), true scale: a single dominant asperity $\sim$17 km SSW of
the centroid nucleation ($\sim$45 km SSW of the epicentre), up-dip, slip at 85--125 km depth, peak 1.63 m,
cleanly contained (edge slip 1--3\,\%). Effective area 2695 km$^2$ $\to$ radius 29 km $\to$
stress drop $\approx$2.9 MPa.}
\end{figure}
\begin{figure}[H]\centering
\includegraphics[width=0.85\textwidth]{../20260810123502/ffm.0/NP1_shift/plots/SlipDist_plane0_5s.png}
\caption{Same model with 5-s rupture-front isochrones: main-phase duration T$_{5-95}$
$\approx$ 18.7 s (times relative to the centroid nucleation; the $\sim$35-s precursor precedes
this and is below our band).}
\end{figure}
\begin{figure}[H]\centering
\includegraphics[width=0.55\textwidth]{../20260810123502/ffm.0/NP1_shift/plots/MomentRate.png}
\caption{Main-phase moment-rate function (centroid-anchored, preferred plane).}
\end{figure}
\begin{figure}[H]\centering
\includegraphics[width=0.9\textwidth]{../20260810123502/ffm.0/NP1_shift/plots/Map.png}
\caption{Map view of the preferred-plane slip.}
\end{figure}
\begin{figure}[H]\centering
\includegraphics[width=0.92\textwidth]{../20260810123502/ffm.0/NP1_shift/plots/P_body_waves.png}
\caption{Teleseismic P fits (preferred plane, shifted): the main packet and its depth phases
are now aligned and fit.}
\end{figure}
\begin{figure}[H]\centering
\includegraphics[width=0.92\textwidth]{../20260810123502/ffm.0/NP1_shift/plots/SH_body_waves.png}
\caption{Teleseismic SH fits (preferred plane, shifted).}
\end{figure}
\begin{figure}[H]\centering
\includegraphics[width=0.85\textwidth]{../20260810123502/ffm.0/NP2_shift/plots/SlipIsochrone_trueScale.png}
\caption{Conjugate plane (124/59/159) for comparison: 8.2\,\% worse fit and 40\,\% of peak
slip at its top edge.}
\end{figure}

\section{Moment tensor (MTUQ): the point source breaks on this complex source}
\begin{table}[H]\centering\small
\begin{tabular}{llll}
\toprule
variant & mechanism & \Mw & Kagan vs W-phase \\
\midrule
joint DC @120.5 km & 202/81/$-$65 & 7.10 (grid edge) & 95.3$^\circ$ \\
body P$+$SH @120.5 km & 220/61/$-$52 & 7.10 (grid edge) & 83.6$^\circ$ \\
body depth scan (60--150 km) & minimum at 70--90 km & 7.05--7.25 & unstable 28.6--93$^\circ$ \\
\bottomrule
\end{tabular}
\caption{The body-wave depth curve bottoms at 70--90 km --- shallower than the W-phase
centroid --- and the mechanism is stable nowhere except $\sim$90 km, where it snaps to 225/71/4
(28.6$^\circ$). A 40-km-offset late asperity, a 35-s precursor and a 40-km-tall slip zone
violate the point-source approximation. The robust element --- strike $\sim$207--225$^\circ$ ---
matches the preferred finite-fault plane. Pipeline rule: for complex sources, take mechanism
and depth from the W-phase and spend the teleseismic data on the finite fault.}
\end{table}

\begin{figure}[H]\centering
\includegraphics[width=0.9\textwidth]{../fig_colombia_mtuq_beachballs.png}
\caption{W-phase reference vs MTUQ attempts.}
\end{figure}
\begin{figure}[H]\centering
\includegraphics[width=\textwidth]{../fig_colombia_assessment.png}
\caption{Left: the main-phase STF is stable across the two assumed planes. Centre: the
point-source depth curve vs the finite-source depth extent --- the point source is biased
shallow by finiteness. Right: the plane-discrimination margin in the pipeline's rupture-size
context.}
\end{figure}

\section{Directivity: two-stage, unilateral southward $+$ up-dip}
(A correction made within the same pre-USGS freeze window: the first push described the
asperity as ``NE along strike'' --- an axis-misread of the WISP plot; the geographic
computation is authoritative.)
\begin{enumerate}
\item \textbf{Model-independent stage 1:} epicentre $\to$ W-phase centroid $=$ 30 km toward
azimuth 163$^\circ$ (SSE), 110 $\to$ 120.5 km, in $\sim$34.5 s --- apparent 0.9 km s$^{-1}$
(the emergent nucleation phase).
\item \textbf{Preferred plane stage 2:} the early$\to$late moment centroid moves 27 km toward
azimuth 220$^\circ$ (SW) while shallowing 110 $\to$ 101 km; 75\,\% of the moment lies forward
of the nucleation vs 24\,\% behind --- predominantly unilateral, along-strike SW plus up-dip.
\item \textbf{Robustness:} the up-dip component appears on both planes; the conjugate's
horizontal direction (northward) contradicts the model-independent stage-1 offset, further
disfavouring that already-rejected plane.
\end{enumerate}

\section{Tectonic context: background seismicity, the Caldas tear, and tomography}
Background: 1{,}981 ComCat events (M$\geq$4, depth $\geq$60 km, 1980--2026) and the CARMA
finite-frequency $dV_p$ model of Cornthwaite et al.\ (2021; IRIS EMC
\texttt{CARMA\_nwSA\_dVp\_Tomo}), the same array behind Sun et al.\ (2022). (The EMC file
labels its units ``percent'' but stores fractions; values are rescaled here --- max $|dV_p|$
$\approx$ 9.9\,\%, matching the publications.)

\begin{figure}[H]\centering
\includegraphics[width=\textwidth]{../fig_colombia_context.png}
\caption{Left: intermediate-depth seismicity; the rupture (red) lies inside the NNE-trending
Cauca cluster, which terminates at the Caldas tear ($\sim$5.15$^\circ$N); the Bucaramanga nest
is the dense knot 407 km NE. Centre: CARMA $dV_p$ at 120 km --- the strong fast anomaly
underlies Bucaramanga (the Caribbean slab of Sun et al.\ 2022); amplitude at our event is
$\sim$0.01\,\% (SW resolution edge). Right: E--W section at 4.84$^\circ$N --- the slip patch
(85--125 km) sits exactly on the east-dipping seismicity band that defines the local Nazca
slab.}
\end{figure}

Four statements follow. (1) \textbf{The rupture is the Cauca cluster breaking}: the slip zone
coincides with the cluster's seismicity band in map view and in depth, with the hypocentre at
its northern/updip edge and rupture propagating SSW \emph{into} the cluster interior, away
from the tear. (2) \textbf{The Caldas tear bounds the system}: the Cauca band terminates at
$\sim$5.15$^\circ$N and a marked gap in deep seismicity (5.2--6.5$^\circ$N) separates it from
the Bucaramanga domain --- consistent with the event's along-strike-compression (map-view
strike-slip) mechanism adjacent to the tear. (3) \textbf{Bucaramanga is a different system}:
per Sun et al.\ (2022) the nest lies in the \emph{Caribbean} slab (the fast anomaly in the
centre panel), so nest and event are plausibly different subducted plates, not merely
different segments. (4) \textbf{Tomographic coverage, honestly}: the CARMA model's domain
nominally contains the event but its amplitude there is at the resolution edge; the local slab
is defined by seismicity, not tomography.

\section{Frozen prospective prediction}
\begin{itemize}
\item \textbf{Fault plane: 226/72/33} (NE--SW, steep, oblique reverse) --- 8.2\,\% margin plus
the conjugate's edge violation.
\item \textbf{Two-stage rupture:} a $\sim$35-s weak emergent nucleation from the cataloged
hypocentre, then the main asperity \textbf{$\sim$45 km SSW of the epicentre, up-dip, 85--125 km depth}; peak slip
$\approx$1.6 m; main-phase T$_{5-95}\approx$19 s; $\Delta\sigma\approx$2.9 MPa.
\item \textbf{$M_0$} $1.68\times10^{20}$ N\,m (\Mw~7.42), matching the W-phase to 1\,\%.
\item \textbf{Scorecard when USGS posts:} plane choice; whether they relocate the
hypocentre/nucleation; asperity position (SSW $+$ up-dip); depth extent; peak slip; duration.
\end{itemize}

\textbf{Caveats.} Teleseismic-only; the precursor's own moment and location are below our band
and not imaged; the mechanism statement rests on the W-phase (MTUQ instability documented
above); body-only fits on both planes (surface-wave GF cap).

\end{document}
